\begin{figure}[htbp]
\centering
\begin{tikzpicture}[
 box/.style={draw=black!70,rounded corners=1pt,align=center,
             font=\small,inner sep=5pt,text width=40mm},
 arrow/.style={-{Stealth[length=2mm]},thick}]
\node[box] (memory) at (0,0) {Charge et mémoires\\\((q,m_0,m_1)\)};
\node[box] (register) at (0,-1.65) {Registre reconstruit\\\(K=(q/12)\mathbf1+Lm\)};
\node[box] (weights) at (-3.8,-3.55) {Partition ordonnée\\\(d_j=K_j/q\)};
\node[box] (direction) at (3.8,-3.55) {Direction de somme nulle\\\(u=P_3P_{11}K\)};
\node[box] (geometry) at (-4.6,-6) {Mesure de frontière\\et forme canonique\\\(C(d),\ \Omega_{P_K}\)};
\node[box] (energy) at (0,-6) {Énergie nodale\\et détails\\\(L_d,\ \mathcal E(z)\)};
\node[box] (lattice) at (4.6,-6) {Déformation réticulaire\\et réseau dual\\\(\Lambda_s^*=\Lambda_{-s}\)};
\draw[arrow] (memory)--(register);
\draw[arrow] (register)--(weights);
\draw[arrow] (register)--(direction);
\draw[arrow] (weights)--(geometry);
\draw[arrow] (weights)--(energy);
\draw[arrow] (direction)--(lattice);
\draw[arrow] (direction) to[bend left=12] node[right,font=\scriptsize,align=left,fill=white,inner sep=1pt]{flux\\hérité} (energy);
\end{tikzpicture}
\caption{Réalisations ultérieures du registre récupérable. La branche d'intervalles détermine les mineurs de chemins, la forme canonique du polygone et le laplacien nodal. La direction algébrique détermine le flot des intervalles et, avec la réalisation réticulaire marquée, la collection duale à vingt-quatre dimensions. Chaque flèche conserve le domaine et les hypothèses de sa construction ; les codomaines inférieurs sont distincts.}
\label{fig:multimorphology}
\end{figure}
